\chapter{El papel de la doble acción integral}\label{chap:doble_integral}

La hipótesis de esta tesis tiene dos partes (capítulo~\ref{chap:introduccion}). La primera, que el PII estabiliza el levitador en presencia de la zona muerta y de la fricción, se confirmó en el capítulo~\ref{chap:desempeno}. La segunda sostiene que la segunda integración reduce, además, el tiempo de atascamiento y la amplitud del ciclo límite que el capítulo~\ref{chap:desempeno} describió. Este capítulo la somete a prueba y, como el resultado la contradice, explica por qué. La pregunta que lo organiza es qué aporta realmente la segunda acción integral.

El recorrido tiene cuatro pasos. La sección~\ref{sec:pi_pii} compara el PII con un controlador PI en las mismas pruebas, y la sección~\ref{sec:barrido_familias} verifica que el resultado no depende de la sintonización particular de cada uno. La sección~\ref{sec:analisis_ciclo} identifica el mecanismo que produce el ciclo límite y que explica por qué la segunda integración no lo modifica. La sección~\ref{sec:ventaja_rampa} señala el único escenario en el que esa integración sí ofrece una ventaja.

\section{Comparación directa: PI frente a PII}\label{sec:pi_pii}

Las pruebas de regulación y seguimiento del capítulo~\ref{chap:desempeno} muestran que el PII estabiliza el sistema, pero no permiten saber qué aporta la segunda integración respecto a un controlador con una sola. Para averiguarlo, ambas pruebas se repitieron con el controlador PI de comparación, definido en la sección~\ref{sec:pi} (ecuación~\eqref{eq:pi}), en las mismas condiciones: escalón de $-2$ a $-3$~cm y un periodo de la referencia trapezoidal, con zona muerta, fricción de Karnopp y $u\in[0;\,3{,}5]$~V. La tabla~\ref{tab:pi_pii} resume los resultados.

\begin{table}[ht]
\centering
\caption{Comparación entre los controladores PI y PII con atascamiento-deslizamiento.}
\label{tab:pi_pii}
\begin{tabular}{|l|c|c|}
\hline
 & PI & PII \\ \hline
Sobrepaso ante el escalón & $25{,}9$\,\% & $23{,}6$\,\% \\ \hline
Ciclo límite pico a pico (escalón) & $0{,}031$~cm & $0{,}033$~cm \\ \hline
Duración media de cada atascamiento (escalón) & $0{,}26$~s & $0{,}36$~s \\ \hline
Duración media de cada atascamiento (trapezoidal) & $0{,}29$~s & $0{,}39$~s \\ \hline
Atascamiento más largo (trapezoidal) & $1{,}07$~s & $3{,}13$~s \\ \hline
Fracción del tiempo atascado (trapezoidal) & $74$\,\% & $76$\,\% \\ \hline
Error eficaz en las rampas & $0{,}0112$~cm & $0{,}0118$~cm \\ \hline
Error medio en las rampas & $-3\times10^{-5}$~cm & $1\times10^{-5}$~cm \\ \hline
\end{tabular}
\end{table}

Los resultados no muestran la ventaja que se esperaba de la segunda acción integral. El PII reduce ligeramente el sobrepaso, pero los atascamientos son en promedio más largos que con el PI y el ciclo límite tiene prácticamente la misma amplitud. Ambos controladores operan sin saturar y con errores eficaces del mismo orden. En las rampas, el error residual que un controlador de tipo uno debe producir es muy inferior a la amplitud del ciclo límite, de modo que con esta referencia los dos son indistinguibles (sección~\ref{sec:ventaja_rampa}).

\section{¿Es un efecto de la sintonización?}\label{sec:barrido_familias}

Un resultado de este tipo puede deberse a que uno de los dos controladores quedó peor sintonizado que el otro. Para descartarlo, se evaluaron dos familias de controladores construidas a partir de los anteriores. En cada una se escaló la ganancia por un factor $k_K$ entre $0{,}5$ y $3$, y se desplazaron los ceros de baja frecuencia, los asociados a la acción integral, por un factor $\alpha$ entre $0{,}25$ y $4$. De las combinaciones resultantes, 256 son estables con un margen de fase de al menos $30^\circ$ en todo el intervalo de $-3$ a $-2$~cm. De ellas, 54 PI y 51 PII completan, sin saturar ni perder la estabilidad, el escalón de $-2$ a $-3$~cm y una rampa de $-0{,}05$~cm/s, doce veces más pronunciada que la de la referencia trapezoidal. El barrido se simuló en una unidad de procesamiento gráfico (GPU), y la validación con integración \texttt{ode45} de los controladores extremos reproduce sus valores con diferencias menores al 8\,\%. La figura~\ref{fig:r15_familias} muestra, para cada controlador, la amplitud del ciclo límite y la duración media de los atascamientos.

\begin{figure}[ht]
    \centering
    % Generado a partir de calculus/Final_Bien/R15_sintonizacion (familias_R15.m, simular_R15.py).
\definecolor{tray}{RGB}{31,94,166}
\definecolor{ctl2}{RGB}{196,96,40}
\begin{tikzpicture}
\begin{axis}[width=0.85\textwidth, height=7cm, xmin=0, xmax=0.12, ymin=0, ymax=1.5,
    xlabel={ciclo límite pico a pico [\si{cm}]}, ylabel={duración media del atascamiento [\si{s}]},
    grid=major, grid style={gray!25}, tick label style={font=\small}, label style={font=\small},
    legend style={font=\small, at={(0.98,0.98)}, anchor=north east}, legend cell align=left,
    /pgf/number format/use comma, scaled ticks=false, x tick label style={/pgf/number format/fixed, /pgf/number format/precision=2},
    table/col sep=comma]
  \addplot[only marks, mark=o, mark size=2.2pt, ctl2] table[x=ciclo, y=atasc] {images/r15_results/tikz/familia_pi.csv};
  \addlegendentry{familia PI}
  \addplot[only marks, mark=square, mark size=2pt, tray] table[x=ciclo, y=atasc] {images/r15_results/tikz/familia_pii.csv};
  \addlegendentry{familia PII}
  \addplot[only marks, mark=*, mark size=3.5pt, ctl2, draw=black] coordinates {(0.0309,0.259)};
  \addlegendentry{PI de la comparación}
  \addplot[only marks, mark=square*, mark size=3.2pt, tray, draw=black] coordinates {(0.0339,0.352)};
  \addlegendentry{PII de la tesis}
\end{axis}
\end{tikzpicture}

    \caption[Familias PI y PII]{Ciclo límite y duración media de los atascamientos en el escalón de $-2$ a $-3$~cm para las familias de controladores PI y PII que operan sin saturar ($u\in[0;\,3{,}5]$~V). Los símbolos rellenos corresponden a los controladores de la tabla~\ref{tab:pi_pii}.}
    \label{fig:r15_familias}
\end{figure}

El ciclo límite más pequeño de la familia PII ($0{,}026$~cm) duplica el del mejor PI ($0{,}013$~cm), y para amplitudes de ciclo límite similares los PI se atascan durante menos tiempo. Ninguna variante del PII mejora, por tanto, a la mejor del PI, y la amplitud del ciclo límite disminuye al aumentar la ganancia proporcional, no al añadir una segunda integración. La diferencia no es un artefacto de la sintonización. Falta explicar por qué ocurre, y a eso se dedica la sección siguiente.

\section{Análisis del ciclo límite}
\label{sec:analisis_ciclo}
% Datos: calculus/Final_Bien/R15_sintonizacion/
% (analisis_ciclo.m, simular_R15.py, exportar_ciclo.py).

La comparación deja dos hechos por explicar. En presencia de atascamiento-deslizamiento, el imán oscila alrededor de la referencia con una amplitud de unas centésimas de centímetro, sin quedar nunca en reposo sobre ella, y la segunda acción integral no reduce esa oscilación respecto a un PI. Esta sección los explica paso a paso. Primero se traduce la fricción estática a términos de voltaje. Después se describe un ciclo completo, se identifica la acción integral como su causa y se estiman la duración de cada atascamiento y la amplitud del ciclo. Todas las cifras corresponden al escalón de $-2$ a $-3$~cm con atascamiento-deslizamiento. Las estadísticas de los ciclos se calculan en el estado estacionario, entre 20 y 40~s, por lo que difieren ligeramente de las de la tabla~\ref{tab:pi_pii}, que abarcan todo el intervalo posterior al escalón.

%-----------------------------------------------------------------------------
\subsection{La fricción estática vista desde el voltaje}

Mientras el imán está atascado, el modelo de Karnopp lo mantiene inmóvil siempre que la fuerza neta $F_e = u/[b(a-x)^4] - mg$ no supere la fricción estática máxima $F_s$. Despejando el voltaje, el imán permanece atascado en la posición $x$ mientras
\begin{equation}
u_e(x)\left(1-\frac{F_s}{mg}\right) \;\le\; u \;\le\; u_e(x)\left(1+\frac{F_s}{mg}\right) .
\label{eq:banda_voltaje}
\end{equation}
El centro de la banda~\eqref{eq:banda_voltaje} es el voltaje de equilibrio $u_e(x)=mgb(a-x)^4$, el que sostiene al imán en la posición $x$. La fricción estática se manifiesta así como una banda muerta de voltaje centrada en $u_e$, con una anchura relativa de $\pm F_s/(mg)=\pm14{,}5$\,\%. En $x_1=-3$~cm, donde $u_e=2{,}393$~V, la banda va de $2{,}047$ a $2{,}739$~V: mientras el voltaje se mantenga dentro de ese intervalo de $0{,}69$~V, el imán no se mueve. Para despegarlo, el controlador tiene que llevar el voltaje hasta uno de los bordes, es decir, cambiarlo al menos $u_e F_s/(mg)=0{,}346$~V respecto al equilibrio.

%-----------------------------------------------------------------------------
\subsection{Anatomía de un ciclo}

La figura~\ref{fig:anatomia} muestra tres segundos del estado estacionario con el PII. Cada ciclo consta de cuatro fases que se repiten de forma casi idéntica.

\begin{enumerate}
  \item \textbf{Atascamiento.} El imán queda detenido a un lado de la referencia, con un error de unos $0{,}016$~cm. La velocidad es nula y el error es constante, pero no nulo, de modo que la acción integral sigue modificando el voltaje: en la gráfica inferior se observa una rampa que se aleja de $u_e$.
  \item \textbf{Ruptura.} Cuando el voltaje alcanza el borde de la banda~\eqref{eq:banda_voltaje}, la fuerza neta supera $F_s$ y el imán se despega. En promedio, el voltaje cambia $0{,}34$~V durante cada atascamiento, prácticamente la mitad de la anchura de la banda.
  \item \textbf{Deslizamiento.} Una vez en movimiento, la fricción pasa de estática a viscosa y desaparece de golpe la fuerza que equilibraba al imán. La fuerza en exceso lo acelera hacia la referencia, y el controlador no puede retirar instantáneamente el voltaje que acumuló durante el atascamiento: el imán cruza la referencia y recorre unos $0{,}032$~cm.
  \item \textbf{Nuevo atascamiento.} Cuando la velocidad vuelve a anularse, el imán queda detenido del otro lado de la referencia y el proceso se repite en sentido contrario.
\end{enumerate}

Cada ciclo completo tiene dos atascamientos y dura en promedio $0{,}98$~s con el PII y $0{,}74$~s con el PI. La amplitud pico a pico ($0{,}033$ y $0{,}031$~cm, respectivamente) coincide con la distancia recorrida en cada deslizamiento.

\begin{figure}[ht]
    \centering
    % Estilo común de las figuras del análisis del ciclo límite (pgfplots).
\pgfplotsset{
  ciclo/.style={
    width=0.9\textwidth, height=4.2cm,
    grid=major, grid style={gray!25},
    tick label style={font=\small}, label style={font=\small}, legend style={font=\small},
    /pgf/number format/use comma, scaled ticks=false,
    y tick label style={/pgf/number format/fixed},
    x tick label style={/pgf/number format/fixed},
    table/col sep=comma, unbounded coords=jump,
    every axis plot/.append style={line join=round},
  },
}
\definecolor{tray}{RGB}{31,94,166}
\definecolor{refc}{RGB}{40,40,40}
\definecolor{ctl2}{RGB}{196,96,40}
\definecolor{ver}{RGB}{46,133,64}

\begin{tikzpicture}
\begin{axis}[ciclo, name=p1, xmin=20, xmax=23, xticklabels={}, ylabel={$x_1$ [\si{cm}]}, ymin=-3.025, ymax=-2.975,
    y tick label style={/pgf/number format/precision=2}]
  \addplot[refc, densely dashed] coordinates {(20,-3) (23,-3)};
  \addplot[tray, line width=1pt] table[x=t, y=x] {images/ciclo_results/tikz/anatomia_pii.csv};
\end{axis}
\begin{axis}[ciclo, name=p2, at={(p1.south west)}, anchor=north west, yshift=-0.25cm, xmin=20, xmax=23, xticklabels={},
    ylabel={$x_2$ [\si{cm/s}]}]
  \addplot[tray, line width=0.8pt] table[x=t, y=v] {images/ciclo_results/tikz/anatomia_pii.csv};
\end{axis}
\begin{axis}[ciclo, at={(p2.south west)}, anchor=north west, yshift=-0.25cm, xmin=20, xmax=23, ymin=1.88, ymax=2.92,
    xlabel={tiempo [\si{s}]}, ylabel={$u$ [\si{V}]}]
  \addplot[ver, densely dashed, line width=0.9pt] coordinates {(20,2.0473) (23,2.0473)};
  \addplot[ver, densely dashed, line width=0.9pt] coordinates {(20,2.7394) (23,2.7394)};
  \addplot[refc!60, densely dotted] coordinates {(20,2.3934) (23,2.3934)};
  \addplot[tray, line width=1pt] table[x=t, y=u] {images/ciclo_results/tikz/anatomia_pii.csv};
  \node[font=\scriptsize, anchor=south east, ver!70!black] at (axis cs:22.98,2.7394) {$u_e(1+F_s/mg)$};
  \node[font=\scriptsize, anchor=north east, ver!70!black] at (axis cs:22.98,2.0473) {$u_e(1-F_s/mg)$};
  \node[font=\scriptsize, anchor=south west, refc!80] at (axis cs:20.02,2.3934) {$u_e$};
\end{axis}
\end{tikzpicture}

    \caption[Anatomía del ciclo límite]{Tres segundos del ciclo límite con el PII en $x_1=-3$~cm. De arriba abajo: posición, velocidad y voltaje. Las líneas verdes discontinuas son los bordes de la banda~\eqref{eq:banda_voltaje}: el imán solo se despega cuando el voltaje llega a uno de ellos.}
    \label{fig:anatomia}
\end{figure}

%-----------------------------------------------------------------------------
\subsection{La acción integral es la causa del ciclo}

El mecanismo anterior depende de que el voltaje siga cambiando mientras el imán está detenido. En un controlador sin acción integral, la salida depende solo del error actual y de su derivada, y ambos permanecen constantes durante el atascamiento. Para comprobarlo, se suprimieron del PII los dos integradores y los dos ceros de baja frecuencia que los acompañan, lo que deja el controlador proporcional-derivativo (PD)
\begin{equation}
C_{PD}(s) = \frac{152{,}43\,(s+17{,}31)}{s+301{,}1} ,
\label{eq:pd}
\end{equation}
que estabiliza el lazo en todo el intervalo de operación. Como un PD no puede generar por sí mismo el voltaje de equilibrio, se le sumó la prealimentación $u_e(r)$ correspondiente a la referencia.

La figura~\ref{fig:pd_pi_pii} compara los tres controladores. Con el PD, el imán llega a la nueva referencia, se atasca y permanece inmóvil, sin ciclo límite. A cambio, el error final puede ser distinto de cero. La ganancia estática del PD, $8{,}76$~V/cm, equivale a una rigidez de $506{,}5$~kg\,cm/s$^2$ por centímetro de error, de modo que el imán puede quedar detenido en cualquier punto donde el error sea menor que $F_s/506{,}5 = 0{,}039$~cm. En esta simulación se detuvo a $0{,}002$~cm de la referencia.

\begin{figure}[ht]
    \centering
    % Estilo común de las figuras del análisis del ciclo límite (pgfplots).
\pgfplotsset{
  ciclo/.style={
    width=0.9\textwidth, height=4.2cm,
    grid=major, grid style={gray!25},
    tick label style={font=\small}, label style={font=\small}, legend style={font=\small},
    /pgf/number format/use comma, scaled ticks=false,
    y tick label style={/pgf/number format/fixed},
    x tick label style={/pgf/number format/fixed},
    table/col sep=comma, unbounded coords=jump,
    every axis plot/.append style={line join=round},
  },
}
\definecolor{tray}{RGB}{31,94,166}
\definecolor{refc}{RGB}{40,40,40}
\definecolor{ctl2}{RGB}{196,96,40}
\definecolor{ver}{RGB}{46,133,64}

\begin{tikzpicture}
\begin{axis}[ciclo, name=p1, height=5cm, xmin=0, xmax=40, ymin=-3.3, ymax=-1.9, xlabel={},
    ylabel={posición [\si{cm}]}, legend style={at={(0.99,0.97)}, anchor=north east}, legend cell align=left, legend columns=3]
  \addplot[ver, line width=1.1pt] table[x=t, y=x_pd] {images/ciclo_results/tikz/pd_pi_pii.csv}; \addlegendentry{PD + $u_e(r)$}
  \addplot[ctl2, line width=0.7pt] table[x=t, y=x_pi] {images/ciclo_results/tikz/pd_pi_pii.csv}; \addlegendentry{PI}
  \addplot[tray, line width=0.7pt] table[x=t, y=x_pii] {images/ciclo_results/tikz/pd_pi_pii.csv}; \addlegendentry{PII}
  \fill[ctl2, opacity=0.12] (axis cs:15,-3.3) rectangle (axis cs:25,-1.9);
\end{axis}
\begin{axis}[ciclo, at={(p1.south west)}, anchor=north west, yshift=-0.45cm, height=4.2cm, xmin=15, xmax=25,
    ymin=-3.025, ymax=-2.975, xlabel={tiempo [\si{s}]}, ylabel={detalle [\si{cm}]}, y tick label style={/pgf/number format/precision=2}]
  \addplot[refc, densely dashed] coordinates {(15,-3) (25,-3)};
  \addplot[ctl2, line width=0.7pt] table[x=t, y=x_pi] {images/ciclo_results/tikz/pd_pi_pii.csv};
  \addplot[tray, line width=0.7pt] table[x=t, y=x_pii] {images/ciclo_results/tikz/pd_pi_pii.csv};
  \addplot[ver, line width=1.4pt] table[x=t, y=x_pd] {images/ciclo_results/tikz/pd_pi_pii.csv};
\end{axis}
\end{tikzpicture}

    \caption[PD, PI y PII con atascamiento-deslizamiento]{Escalón de $-2$ a $-3$~cm con el PD con prealimentación, el PI y el PII. Abajo, detalle de 15 a 25~s: el PD queda en reposo y los controladores con acción integral oscilan alrededor de la referencia.}
    \label{fig:pd_pi_pii}
\end{figure}

El ECP-730 reproduce así el \emph{hunting} que la sección~\ref{sec:fund_friccion} describe en general. Sin acción integral, el lazo se detiene con un error que depende de dónde quedó atascado y de qué tan exacto sea el voltaje de equilibrio que se suministra. Con acción integral, el error medio se anula aunque el modelo no sea exacto, pero el lazo nunca se detiene, porque el integrador sigue acumulando el error residual hasta forzar una nueva ruptura. El control conmutado de Hernández Alcántara et al.~\cite{alcantara} aborda precisamente este dilema, al desactivar la acción integral cuando el error es pequeño.

%-----------------------------------------------------------------------------
\subsection{Duración del atascamiento}

Si el error con el que el imán se atasca es $e_0$, la duración del atascamiento es el tiempo que tarda la acción integral en cambiar el voltaje los $0{,}346$~V necesarios para la ruptura. A bajas frecuencias, el PI de la ecuación~\eqref{eq:pi} se comporta como $K_i/s$ con $K_i=86{,}2$~V/(cm\,s), y el PII como $K_{ii}/s^2 + K_i/s$ con $K_{ii}=79{,}5$~V/(cm\,s$^2$) y $K_i=63{,}4$~V/(cm\,s). Con el error constante, el cambio de voltaje durante el atascamiento es
\begin{equation}
\Delta u_{PI}(t) = K_i e_0\, t, \qquad
\Delta u_{PII}(t) = K_i e_0\, t + \frac{K_{ii} e_0}{2}\, t^2 .
\label{eq:rampa_atasco}
\end{equation}

Para el PI, con $e_0=0{,}0151$~cm, el voltaje crece a $K_i e_0 = 1{,}30$~V/s y alcanza el borde de la banda, que en promedio exige un cambio de $0{,}328$~V, en $0{,}328/1{,}30=0{,}25$~s, en acuerdo con la duración media medida de $0{,}26$~s. Para el PII, la figura~\ref{fig:atasco} muestra que el voltaje sigue la parábola del cambio de voltaje del PII~\eqref{eq:rampa_atasco}: la curvatura medida coincide con $K_{ii}e_0$ dentro de un 10\,\%. Sin embargo, el término cuadrático solo domina sobre el lineal cuando $t > 2K_i/K_{ii} = 1{,}6$~s, y los atascamientos duran unos $0{,}35$~s. En ese intervalo, la segunda integración aporta aproximadamente una quinta parte del cambio de voltaje, y el término lineal del PII es más lento que el del PI porque su $K_i$ es menor (63 frente a 86~V/(cm\,s)). Además, la pendiente inicial medida es un 20 a 25\,\% menor que $K_i e_0$, lo que indica que el estado del controlador al inicio de cada atascamiento conserva parte de la dinámica del deslizamiento previo. El resultado neto es que el PII tarda más que el PI en despegar el imán.

\begin{figure}[ht]
    \centering
    % Estilo común de las figuras del análisis del ciclo límite (pgfplots).
\pgfplotsset{
  ciclo/.style={
    width=0.9\textwidth, height=4.2cm,
    grid=major, grid style={gray!25},
    tick label style={font=\small}, label style={font=\small}, legend style={font=\small},
    /pgf/number format/use comma, scaled ticks=false,
    y tick label style={/pgf/number format/fixed},
    x tick label style={/pgf/number format/fixed},
    table/col sep=comma, unbounded coords=jump,
    every axis plot/.append style={line join=round},
  },
}
\definecolor{tray}{RGB}{31,94,166}
\definecolor{refc}{RGB}{40,40,40}
\definecolor{ctl2}{RGB}{196,96,40}
\definecolor{ver}{RGB}{46,133,64}

\begin{tikzpicture}
\begin{axis}[ciclo, height=6.2cm, xmin=0, xmax=0.45, ymin=0, ymax=0.42,
    xlabel={tiempo desde el inicio del atascamiento [\si{s}]}, ylabel={cambio de voltaje [\si{V}]},
    legend style={at={(0.98,0.03)}, anchor=south east}, legend cell align=left]
  \addplot[ver, densely dashed, line width=0.9pt] coordinates {(0,0.3461) (0.45,0.3461)};
  \addlegendentry{$u_e F_s/(mg)=0{,}346$ V}
  \addplot[ctl2, line width=0.8pt] table[x=tau, y=du] {images/ciclo_results/tikz/atasco_pi.csv};
  \addlegendentry{PI (simulación)}
  \addplot[ctl2!70!black, dashed, line width=1pt] table[x=tau, y=du] {images/ciclo_results/tikz/atasco_pi_teo.csv};
  \addlegendentry{PI: $K_i e_0 t$}
  \addplot[tray, line width=0.8pt] table[x=tau, y=du] {images/ciclo_results/tikz/atasco_pii.csv};
  \addlegendentry{PII (simulación)}
  \addplot[tray!70!black, dashed, line width=1pt] table[x=tau, y=du] {images/ciclo_results/tikz/atasco_pii_teo.csv};
  \addlegendentry{PII: $K_i e_0 t + K_{ii} e_0 t^2/2$}
\end{axis}
\end{tikzpicture}

    \caption[Voltaje durante el atascamiento]{Cambio de voltaje durante cada atascamiento, alineado al instante en que comienza y con el signo del error. Las curvas discontinuas son la predicción~\eqref{eq:rampa_atasco} con el error medio de atascamiento. La ruptura ocurre al alcanzar $u_eF_s/(mg)$.}
    \label{fig:atasco}
\end{figure}

%-----------------------------------------------------------------------------
\subsection{Amplitud del ciclo}

La amplitud del ciclo la determina la fase de deslizamiento. En el instante de la ruptura, el voltaje está en el borde de la banda, de modo que la fuerza neta sobre el imán vale aproximadamente $F_s$ y deja de estar compensada por la fricción. El imán se desplaza hasta que el controlador, a través de su respuesta a la posición, retira esa fuerza en exceso. Cuanto más rígido es el lazo, es decir, cuanto mayor es la ganancia que traduce el error de posición en voltaje, menor es el desplazamiento necesario. Este razonamiento predice que la amplitud del ciclo es aproximadamente inversamente proporcional a la ganancia del controlador, e independiente del número de integradores.

La figura~\ref{fig:ganancia} lo confirma con las familias de controladores de la sección~\ref{sec:barrido_familias} (figura~\ref{fig:r15_familias}), manteniendo fijos los ceros de baja frecuencia ($\alpha=1$) y variando solo el factor de ganancia $k_K$. La amplitud del ciclo decrece aproximadamente como $k_K^{-1{,}2}$ en la familia PI y como $k_K^{-1{,}1}$ en la familia PII. Con la misma ganancia, ambas estructuras producen prácticamente el mismo ciclo límite: $0{,}072$ frente a $0{,}074$~cm con $k_K=0{,}5$, y $0{,}031$ frente a $0{,}034$~cm con $k_K=1$ (valores del barrido en la GPU, que difieren en menos de 3\,\% de los obtenidos con \texttt{ode45}). La ganancia fija la amplitud del ciclo; la segunda integración no la modifica.

\begin{figure}[ht]
    \centering
    % Estilo común de las figuras del análisis del ciclo límite (pgfplots).
\pgfplotsset{
  ciclo/.style={
    width=0.9\textwidth, height=4.2cm,
    grid=major, grid style={gray!25},
    tick label style={font=\small}, label style={font=\small}, legend style={font=\small},
    /pgf/number format/use comma, scaled ticks=false,
    y tick label style={/pgf/number format/fixed},
    x tick label style={/pgf/number format/fixed},
    table/col sep=comma, unbounded coords=jump,
    every axis plot/.append style={line join=round},
  },
}
\definecolor{tray}{RGB}{31,94,166}
\definecolor{refc}{RGB}{40,40,40}
\definecolor{ctl2}{RGB}{196,96,40}
\definecolor{ver}{RGB}{46,133,64}

\begin{tikzpicture}
\begin{loglogaxis}[ciclo, height=6cm, width=0.75\textwidth, xmin=0.45, xmax=1.1, ymin=0.025, ymax=0.085,
    xlabel={factor de ganancia $k_K$}, ylabel={ciclo límite pico a pico [\si{cm}]},
    xtick={0.5,0.6,0.7,0.8,0.9,1}, ytick={0.03,0.04,0.05,0.06,0.08}, log ticks with fixed point,
    legend style={at={(0.98,0.98)}, anchor=north east}, legend cell align=left]
  \addplot[ctl2, mark=o, mark size=2.2pt, line width=0.8pt] table[x=kK, y=ciclo] {images/ciclo_results/tikz/ganancia_pi.csv}; \addlegendentry{familia PI ($\alpha=1$)}
  \addplot[tray, mark=square, mark size=2pt, line width=0.8pt] table[x=kK, y=ciclo] {images/ciclo_results/tikz/ganancia_pii.csv}; \addlegendentry{familia PII ($\alpha=1$)}
  \addplot[refc!60, densely dashed, domain=0.45:1.1, samples=2] {0.032/x}; \addlegendentry{$\propto 1/k_K$}
\end{loglogaxis}
\end{tikzpicture}

    \caption[Ciclo límite contra la ganancia]{Amplitud del ciclo límite contra el factor de ganancia $k_K$, con los ceros de baja frecuencia fijos. Ganancias mayores saturan el actuador.}
    \label{fig:ganancia}
\end{figure}

%=============================================================================
\section{Dónde sí ayuda la segunda integración}\label{sec:ventaja_rampa}

La ventaja teórica del PII aparece al seguir una rampa, como se vio en la sección~\ref{sec:fund_tipo} al estudiar el tipo del sistema. El lazo con el PI es de tipo uno. Ante una rampa de pendiente $\dot r$ conserva un error estacionario $\dot r/K_v$, con la constante de velocidad~\eqref{eq:constantes_error} igual a $K_v=91{,}6$~s$^{-1}$ en $x_1^*=-3$~cm. El lazo con el PII es de tipo dos y anula ese error. Con la pendiente de la referencia trapezoidal ($0{,}004$~cm/s), el error del PI sería de solo $4\times10^{-5}$~cm, muy inferior a la amplitud del ciclo límite, y por eso la tabla~\ref{tab:pi_pii} no distingue entre los dos controladores.

Con la rampa doce veces más pronunciada del barrido ($0{,}05$~cm/s), la predicción para el PI es $5{,}5\times10^{-4}$~cm, y la simulación da un error medio de $5{,}6\times10^{-4}$~cm con el PI frente a $1{,}3\times10^{-4}$~cm con el PII. La segunda integración cumple su función, pero el ciclo límite impone en ambos casos un error eficaz de unos $7{,}6\times10^{-3}$~cm, prácticamente igual con los dos controladores y más de diez veces mayor que esos errores medios. La ventaja del PII no se traduce, por tanto, en un mejor seguimiento. En este sistema solo sería apreciable con referencias de pendiente mucho mayor o con una fricción estática mucho menor.

\section{Síntesis}\label{sec:sintesis_doble_integral}

La segunda parte de la hipótesis se rechaza. La fricción estática equivale a una banda muerta de voltaje de $\pm14{,}5$\,\% alrededor del voltaje de equilibrio, y el imán solo se despega cuando el voltaje alcanza uno de sus bordes. La acción integral es la causa del ciclo límite, porque obliga al lazo a romper esa banda una y otra vez, y la duración de cada atascamiento la fija su término simple, ya que el término doble solo dominaría en atascamientos de más de $1{,}6$~s, y los observados duran unos $0{,}35$~s. La amplitud del ciclo la fija la ganancia del controlador y no el número de integradores. La segunda integración elimina el error en rampa, pero ese error queda muy por debajo del que impone el ciclo límite.

La doble acción integral estabiliza el sistema, pero no es la que determina la calidad de la regulación en presencia de fricción estática. El capítulo~\ref{chap:conclusiones} contrasta esta conclusión con los antecedentes y propone cómo abordar el ciclo límite.

\FloatBarrier
